\documentclass{article}
\usepackage[preprint]{neurips_2023}
\usepackage[utf8]{inputenc}
\usepackage{amsmath}
\usepackage{graphicx}
\usepackage{natbib}

\title{Grokking Cliffs: How Model Collapse Prevents Delayed Generalization}
\author{%
  Research Team \\
  Department of Computer Science \\
}
\date{\today}

\begin{document}

\maketitle

\begin{abstract}
We study how model collapse prevents grokking. Through extensive experiments on modular arithmetic, we demonstrate that there are pure groks at step 1400, while low collapse delays this to step 3100. Severe collapse entirely prevents grokking. We observe corresponding weight norm drops of 30-42\% during the generalization phase. We present a theoretical framework mapping collapse severity to grokking probability and predicting the exact phase transition cliff.
\end{abstract}

\section{Introduction}
Grokking, the phenomenon of delayed generalization long after overfitting, has been primarily studied in pristine data environments. However, as generative models increasingly train on synthetic data, understanding the intersection of model collapse and grokking is crucial.

\section{Related Work}
Previous works have established the foundational dynamics of grokking \citep{power2022grokking}. Concurrently, the phenomenon of model collapse resulting from recursive training on AI-generated data has been formalized \citep{shumailov2024curse}. This work bridges these domains.

\section{Method}
We trained a 1-layer transformer (214K parameters) on a modular arithmetic task under varying levels of label noise and synthetic contamination. We track test accuracy, Fourier concentration, and weight norms.

\section{Theory}
We introduce a dynamical systems model characterizing the weight norm trajectory $w(t)$ as a differential equation competing between weight decay $\lambda$ and noise gradients $\eta$. The mutual information between the weights and the true task degrades as a function of the collapse severity.
\begin{equation}
\frac{dw}{dt} = -\lambda w + \eta
\end{equation}

\section{Results}
Our results show a sharp phase transition:
\begin{itemize}
    \item Pure data: Groks at step 1400.
    \item Low collapse (5\%): Groks at step 3100 (93\% accuracy).
    \item High collapse ($\ge$ 15\%): Never groks. Weight norm remains elevated.
\end{itemize}

\section{Discussion}
Model collapse is not merely a reduction in effective sample size, as scarcity baselines still grok. It acts as an active noise injection that prevents the weight norm from dropping below the critical threshold required for the generalization circuit to dominate.

\bibliographystyle{plainnat}
\bibliography{refs}

\end{document}
